\section {The crescent moon (Astronomy)}

Prove, in the simplest way you can, that the crescent of the moon is not caused by the shadow of the earth falling on the moon
(with the earth standing between the sun and the moon).

\subsection{Solution}

Sometimes, around sunset, you can see both the sun and a crescent moon in the sky at the same time.
You are standing on the earth and see both of them in front of you, so the earth is certainly not between the sun and the moon,
and still the moon is a crescent.

In fact a crescent moon is always close to the sun in the sky: it shows up only shortly after sunset or shortly before sunrise.
The crescent is the half of the moon that faces the sun, seen mostly from the side.
The shadow of the earth does fall on the moon sometimes, but that is a lunar eclipse, which can only happen at full moon,
when the moon is on the opposite side of the sky from the sun.
